% Corte mecánico íntegro: parte_i.tex:420--528 + parte_ii.tex:1--321
% (la línea 322 del propietario es un separador vacío).
% Residencia material del ensamblaje sucesor de 102 capítulos.
% Residencia causal: capítulo 8; no sustituye al propietario Γ_9.
\chapter{TPK: selection, transport, update, and return}

\section{The typed transition}

The TPK is not a nominal container. The phase function and its tabulated vector are, respectively,
\[
 M_{\ph}:\mathcal D_9\longrightarrow\{+1,-1,0\},
 \qquad
 m_{\ph}:=\bigl(M_{\ph}(1),\ldots,M_{\ph}(9)\bigr)^{\mathsf T}.
\]
They are not the same type of object and neither acts by direct composition on the state. For \(x\in X_{\TPK}^{\enr}\), we define
\[
 \operatorname{Tra}_t(y,d):=\operatorname{Lift}(T_d)(y),
 \qquad
 \operatorname{Upd}_t:=\operatorname{Rec}_t\circ\operatorname{Mem}_t.
\]
The elementary transition is then written
\begin{equation}
 U_t(x)=\operatorname{Upd}_t\!\left(
 \operatorname{Tra}_t\!\left(
 \operatorname{Sel}_t\!\left(x,M_{\ph}(g(t))\right),d(t)
 \right)\right),
 \qquad
 U_t:X_{\TPK}^{\enr}\longrightarrow X_{\TPK}^{\enr}.
 \label{espiral:eq:transicion-tpk-tipada}
\end{equation}
The selector \(\operatorname{Sel}_t\) reads the scalar value of the tritic regime and chooses the additive, multiplicative or threshold sheet; \(\operatorname{Tra}_t(\,\cdot\,,d(t))=\operatorname{Lift}(T_{d(t)})\) applies the cardinal direction; \(\operatorname{Mem}_t\) updates residue, quotient, carry, signature, cylinder, boundary and memory; and \(\operatorname{Rec}_t\) performs the typed return. The abbreviated notation \(\mathrm{Upd}_t\circ\mathrm{Tra}_t\circ\mathrm{Sel}_t\) designates this same composition and not a second engine.

The minimal enriched fiber can be represented as
\[
 \mathcal F_q=O_q\times H_q\times C_q\times S_q\times B_q\times\Lambda_q,
\]
with orientation/leaf, lifting, carry, signature, boundary, and record. The
detail of each factor depends on its residence, but the integrity rule
is stable: a coordinate cannot be removed if it changes a continuation.

\section{Nonadic holonomy}

For a nine-step compatible block,
\begin{equation}
 \Gamma_{9,K}=U_{K+8}\circ\cdots\circ U_K.
 \label{espiral:eq:gamma9-def}
\end{equation}
The phase satisfies
\[
 g(K+9)=g(K),
\]
but the enriched state does not in general return:
\[
 \Gamma_{9,K}(x_K)\neq x_K.
\]
In particular,
\[
 q_{\mem}(\Gamma_9x)=q_{\mem}(x)+1,
 \qquad B_{K+9}\neq B_K\quad\text{in general}.
\]

\begin{theorem}[Phase return with memory advance]
The holonomy \(\Gamma_9\) closes the phase coordinate and transports the
state to a new depth. Therefore, phase return is not state
return.
\end{theorem}
\begin{proof}
Phase closure is the periodicity of \(g\). The ledger update
includes the depth increment and boundary transport. If the
state returned, all coordinates would have to coincide, contradicting
\(q_{\mem}\mapsto q_{\mem}+1\). Therefore, only the phase returns.
\end{proof}

\section{Composition of carry}

Let \(\mathcal K_{\Carr}\) be the carry module, and let
\(H,Y:\mathscr P_{\TPK}^{\enr}\to\operatorname{End}(\mathcal K_{\Carr})\)
be the right and left actions induced, respectively, by the preserved prefix
and by the active sheet of a route. We regard
\(\mathcal K_{\Carr}\) as a bimodule and write
\(kH(\gamma)\) and \(Y(\gamma)k\) for those actions. The map
\(\Carr:\mathscr P_{\TPK}^{\enr}\to\mathcal K_{\Carr}\) assigns to each route
its total carry.

For paths composable \(\gamma_1,\gamma_2\), the carry satisfies
\begin{equation}
 \Carr(\gamma_2\gamma_1)
 =\Carr(\gamma_2)H(\gamma_1)
 +Y(\gamma_2)\Carr(\gamma_1).
 \label{espiral:eq:cociclo-carry}
\end{equation}
This law shows what preserving memory means: not keeping a quantity constant, but retaining its exact update rule. Accumulated memory is dynamic and traceable.

\section{Dodecaphase calendar and non-splitting}

The internal calendar is organized through
\[
 0\longrightarrow C_{12}\longrightarrow C_{108}
 \longrightarrow C_9\longrightarrow0.
\]
The extension preserves the difference between a visible turn and the lifted
state. In the realization
\(C_{108}\simeq C_{27}\times C_4\), the part \(C_{27}\to C_9\) does not
split: the lifts of the generator have order \(27\), not an order
dividing \(9\). The factor \(C_4\) supplies the quarter-turn; the ternary
kernel supplies the elevation. The combination is not an external temporal label,
but the return architecture of the TPK.

\section{Posteriority of the Markov quotient as a publication of reduced memory}

Averaging enriched dynamics may produce a reduced-memory operator
\(K_{\ph}\). This operator is useful for studying transfer, stationarity,
and spectrum, but does not generate \(\Gamma_9\). The dependence is
\[
 U_t\longrightarrow\Gamma_9\longrightarrow K_{\ph}.
\]
Inverting it would era sheet, prefix, boundary, and carry. In this publication the
radial reader is constructed over enriched histories, never over the
Markov quotient.

\part{Exact helical geometry of the nonadic memory}

\chapter{Nonadic helix of phase and depth}

\section{Nonadic phase lift to the phase--memory covering}

Let
\[
 \operatorname{Rep}_9=\{0,1,\ldots,8\},
 \qquad \widetilde\Theta=\operatorname{Rep}_9\times\Z.
\]
We define
\begin{equation}
 \widetilde\Gamma_9(r,n)=(r,n+1),
 \qquad
 \Pi(r,n)=[r+9n]_{108}.
 \label{espiral:eq:cobertura-helice}
\end{equation}
The first coordinate records the visible phase and the second records the depth.
It holds
\[
 \Pi(r,n+12)=\Pi(r,n),
 \qquad
 \widetilde\Gamma_9^{12}(r,n)=(r,n+12).
\]
The covering transformation
\[
 \tau_{12}(r,n)=(r,n+12)
\]
acts freely. Each fiber of \(\Pi\) is an orbit of that action.

\begin{definition}[Nonadic helix]
The nonadic helix is the lifted system
\[
 (\widetilde\Theta,\widetilde\Gamma_9,\Pi)
\]
with phase \(r\bmod9\), calendar \(\Pi(r,n)\), and deep memory \(n\).
\end{definition}

\begin{theorem}[Periodic phase and aperiodic memory]
The projection of the helix is periodic of order \(12\) on the calendar \(C_{108}\), while the lifted state distinguishes all depths \(n+12k\), \(k\in\Z\).
\end{theorem}
\begin{proof}
The equality \(\Pi(r,n)=\Pi(r',n')\) forces \(r=r'\) and
\(n-n'\in12\Z\). Therefore, the projection identifies exactly the
orbits of \(\tau_{12}\), but the second coordinate of the lift
remains integral and distinguishes their elements.
\end{proof}

\section{Normalized geometric immersion}

For \(N\ge1\), let
\[
 z_N=q_9(N)+\frac{\rho_9(N)-1}{9}.
\]
The immersion
\begin{equation}
 H_9(N)=\bigl(\cos(2\pi z_N),\sin(2\pi z_N),z_N\bigr)
 \label{espiral:eq:inmersion-helice-nonadica}
\end{equation}
satisfies
\[
 z_{N+9}=z_N+1,
 \qquad
 H_9(N+9)=H_9(N)+(0,0,1).
\]
The projection onto the first two coordinates returns; the third priservis
the turn. The figure is not imposed on the arithmetic: it is an immersion of the
exact decomposition \(N=9q_9(N)+\rho_9(N)\).

For the continuous curve
\[
 h(z)=(\cos2\pi z,\sin2\pi z,z),
\]
we have
\[
 h'(z)=(-2\pi\sin2\pi z,2\pi\cos2\pi z,1),
\]
\[
 h''(z)=(-4\pi^2\cos2\pi z,-4\pi^2\sin2\pi z,0).
\]
Its curvature and torsion are constants:
\begin{equation}
 \kappa_h=\frac{4\pi^2}{4\pi^2+1},
 \qquad
 \mathcal T_h=\frac{2\pi}{4\pi^2+1}.
 \label{espiral:eq:frenet-helice-normalizada}
\end{equation}

\begin{proof}[Calculation of \cref{espiral:eq:frenet-helice-normalizada}]
The norm of \(h'\) is \(\sqrt{4\pi^2+1}\) and
\(\|h'\times h''\|=4\pi^2\sqrt{4\pi^2+1}\). Therefore,
\[
 \kappa_h=\frac{\|h'\times h''\|}{\|h'\|^3}
 =\frac{4\pi^2}{4\pi^2+1}.
\]
Furthermore, with \(h'''=(8\pi^3\sin2\pi z,-8\pi^3\cos2\pi z,0)\),
\[
 \mathcal T_h=
 \frac{\det(h',h'',h''')}{\|h'\times h''\|^2}
 =\frac{2\pi}{4\pi^2+1}.
\]
\end{proof}

\begin{figure}[htbp]
\centering
\begin{tikzpicture}[x={(1cm,0cm)},y={(.30cm,.12cm)},z={(0cm,.68cm)},scale=.82]
 \draw[->,HMTGray] (0,0,0)--(0,0,7.2) node[above] {depth};
 \foreach \k in {0,...,4}{
   \draw[HMTLine] (0,0,{1.4*\k}) ellipse (1.55 and .48);
 }
 \draw[HMTBlue,very thick,domain=0:1440,samples=260,smooth,variable=\t]
   plot ({1.55*cos(\t)},{1.55*sin(\t)},{\t/210});
 \foreach \j in {0,...,35}{
   \fill[HMTRed] ({1.55*cos(40*\j)},{1.55*sin(40*\j)},{40*\j/210}) circle (1.1pt);
 }
 \node[HMTGray,font=\small\sffamily,right] at (1.7,0,3.3) {$9$ marks by turn};
\end{tikzpicture}
\caption{Nonadic helix: the rings are level sections; the curve conserves memory advance.}
\label{espiral:fig:helice-nonadica-analitica}
\end{figure}

\section{Nested circles and cylinders}

A section \(n=\text{constant}\) produces a circle. A family of sections with constant radius produces a cylinder. If the radius depends on depth, the same structure produces a radial helical suspension. Therefore:
\[
 \text{orbit}\longleftrightarrow\text{helical sheet},
\]
\[
 \text{level section}\longleftrightarrow\text{circle},
\qquad
 \text{family of sections}\longleftrightarrow\text{cylinder}.
\]
This is a geometric translation of the lifted system, not a replacement
of the ledger by a drawing.

\chapter{Dynamic suspension of the enriched affine system}

\section{State-dependent holonomy}

Let \(U_{\APP}\) be a radial module and associate with each nonadic block an
affine endomorphism
\[
 H_9(x):u\longmapsto\Lambda_9(x)u+C_9(x).
\]
When holonomy depends on the initial state, we define
\(F(x,u)=(\Gamma_9x,H_9(x)u)\). If \(F\) is continuous, its suspension is
the mapping torus
\begin{equation}
 \mathfrak S_{\TPK}^{\rad}
 =\frac{X_{\TPK}^{\enr}\times U_{\APP}\times[0,1]}
 {(x,u,1)\sim(\Gamma_9x,H_9(x)u,0)}.
 \label{espiral:eq:suspension-skew}
\end{equation}
If \(F\) is to homeomorphism ---which requiris, in addition to the invertibility
of the linear parts, the continuous reversibility of \(\Gamma_9\) and of its
dependence on \(x\)--- to bidirectional suspension is obtained. In the
general monoid the same quotient encodis only forward evolution;
it is not identified with to teliscope and not global inversion is prisupposed.

\begin{definition}[Radial helical suspension]
The suspension \cref{espiral:eq:suspension-skew}, endowed with an angular coordinate
of phase and a subsequent radial evaluation, is called a \emph{radial helical
suspension}. The term \emph{general helix} is reserved for the
cases that satisfy a differential condition of Lancret.
\end{definition}

\section{Orbital sheets and seam of the suspension space}

The orbit of the skew--product is
\[
 (x_0,u_0),
 (\Gamma_9x_0,H_9(x_0)u_0),
 (\Gamma_9^2x_0,H_9(\Gamma_9x_0)H_9(x_0)u_0),\ldots.
\]
The stitching preserves the order of the factors. Replacing the whole orbit with the
final product would prevent truncating the history and reconstructing its prefixes. The
helicoidal sheet is, therefore, an orbit with ledger, not an isolated curve.

\section{Stationary affine case}

If \(H_9(x)=H=(\Lambda,C)\) is constant,
\[
 u_{n+1}=\Lambda u_n+C.
\]
The exact internal iteration, valid for every endomorphism \(\Lambda\), is
\begin{equation}
 u_n=\Lambda^n u_0+\sum_{k=0}^{n-1}\Lambda^kC.
 \label{espiral:eq:orbita-afin-estacionaria-general}
\end{equation}
If \(I-\Lambda\) is invertible in the radial module, the fixed point exists
\[
 u_*=(I-\Lambda)^{-1}C,
\]
and
\begin{equation}
 u_n-u_*=\Lambda^n(u_0-u_*).
 \label{espiral:eq:orbita-afin-estacionaria}
\end{equation}
In the translational sector \(\Lambda=I\),
\[
 u_n=u_0+nC.
\]
Only after choosing a positive scalar character
\(\chi:U_{\APP}\to\R\), with
\(\chi\Lambda=\lambda\chi\), are the numerical formulas
\(u_*=C/(1-\lambda)\) for \(\lambda\neq1\) and the real powers of
\(\lambda\) used in continuous interpolation obtained. An
abstract endomorphism is not divided by \(1-\Lambda\).
These two possibilities anticipate the separation between a
self-scaling family and a translational family. Neither by itself selects a
known spiral; the curve appears after applying \(\exp_p\) and the
angular coordinate.

\begin{figure}[htbp]
\centering
\begin{tikzpicture}[>=Latex,node distance=9mm and 14mm,
 s/.style={draw=HMTNavy,fill=white,minimum width=21mm,minimum height=9mm,font=\small\sffamily}]
\node[s] (x0) {\((x_0,u_0)\)};
\node[s,right=of x0] (x1) {\((x_1,u_1)\)};
\node[s,right=of x1] (x2) {\((x_2,u_2)\)};
\node[right=7mm of x2] (dots) {$\cdots$};
\draw[->,HMTBlue,thick] (x0)--node[above,font=\scriptsize] {$\Gamma_9,H_9(x_0)$} (x1);
\draw[->,HMTBlue,thick] (x1)--node[above,font=\scriptsize] {$\Gamma_9,H_9(x_1)$} (x2);
\draw[->,HMTBlue,thick] (x2)--(dots);
\draw[HMTRed,dashed,rounded corners] ([xshift=-3mm,yshift=-4mm]x0.south west) rectangle ([xshift=3mm,yshift=4mm]x2.north east);
\node[below=7mm of x1,HMTGray,font=\small\sffamily] {each prefix remains accessible};
\end{tikzpicture}
\caption{The suspension preserves the sequence of factors, not only the final holonomy.}
\label{espiral:fig:skew-ledger}
\end{figure}

\chapter{Bidirectional Involution and Conjugate Leaves}

\section{Route inversion}

Let \(\mathscr P_{\TPK}^{\times}\) be the subgroupoid of reversible routes. If
the radial representation satisfies
\[
 h_{e^\dagger}=h_e^{-1},
\]
then, for \(\gamma=e_n\cdots e_1\),
\[
 H_{C_{\mathrm{rev}}\gamma}
 =h_{e_1}^{-1}\cdots h_{e_n}^{-1}
 =(h_{e_n}\cdots h_{e_1})^{-1}
 =H_\gamma^{-1}.
\]

\begin{theorem}[Conjugation of Ascent and Descent]
In the reversible sector, the route involution interchanges the ascent
and descent leaves and inverts the affine holonomy.
\end{theorem}

On the helical coverage, let
\[
 \iota(r,n)=(r,-n).
\]
Then
\begin{equation}
 \iota\widetilde\Gamma_9\iota
 =\widetilde\Gamma_9^{-1}.
 \label{espiral:eq:involucion-helice}
\end{equation}
The equality is obtained directly:
\[
 (r,n)\xmapsto{\iota}(r,-n)
 \xmapsto{\widetilde\Gamma_9}(r,-n+1)
 \xmapsto{\iota}(r,n-1).
\]

\begin{figure}[htbp]
\centering
\begin{tikzpicture}[x={(1cm,0cm)},y={(.30cm,.12cm)},z={(0cm,.58cm)},scale=.72]
 \draw[->,HMTGray] (0,0,-4.8)--(0,0,4.8) node[above] {memory};
 \draw[HMTBlue,very thick,domain=0:900,samples=180,smooth,variable=\t]
   plot ({1.25*cos(\t)},{1.25*sin(\t)},{\t/190});
 \draw[HMTRed,very thick,domain=0:900,samples=180,smooth,variable=\t]
   plot ({1.25*cos(-\t)},{1.25*sin(-\t)},{-\t/190});
 \fill[HMTNavy] (1.25,0,0) circle (1.5pt);
 \node[right,HMTBlue,font=\small\sffamily] at (1.5,0,3.4) {ascent};
 \node[right,HMTRed,font=\small\sffamily] at (1.5,0,-3.4) {descent};
\end{tikzpicture}
\caption{Conjugate leaves produced by route inversion.}
\label{espiral:fig:hojas-conjugadas}
\end{figure}

\section{Mirror of constants and fixed axis}

The involution on the constant fiber
\[
 \mathfrak c(u^\pi,u^e,u^\varphi)
 =(u^e,u^\pi,u^\varphi)
\]
interchanges the branches \(\pi/e\) and fixes the axis \(\varphi\). It also preserves \(u^\pi+u^e\). This involution and route inversion are distinct objects: the former acts on a publication fiber; the latter acts on the path groupoid. No global conjugation \(\mathfrak c\Gamma_9\mathfrak c=\Gamma_9^{-1}\) is asserted without constructing the corresponding intertwiner.

This distinction protects the positive thesis. The mirror \(\pi/e\), the axis
\(\varphi\), and bidirectionality are present; what is not permitted is
to identify them through an untyped equality of operators.

\section{Nested circles as cross sections of the leaves}

Let \(r_n>0\) be the scale of section \(n\). The family
\[
 C_n=\{(r_n\cos\theta,r_n\sin\theta,n):0\le\theta<2\pi\}
\]
produces nested circles as \(r_n\) varies and cylinders when
sections with a common radius are grouped. The union of the \(C_n\) along a TPK
route is the discrete trace of a helical sheet. The involution
\(n\mapsto-n\) produces the conjugate sheet.
